\begin{table}
\centering
\begin{talltblr}[         %% tabularray outer open
caption={Effect of Segregation on County Fiscal Policy Mechanisms: OLS and IV},
note{}={+ p < 0.1, * p < 0.05, ** p < 0.01, *** p < 0.001},
note{ }={OLS and IV (RDI instrument) estimates of the effect of segregation (D) on each county fiscal policy mechanism. Each column is a separate outcome, and the columns are the same policy channels controlled for in the longevity models. All models include decade (1980, 1990, 2000), urban-rural, and region fixed effects. Standard errors clustered at the county level in parentheses. D is scaled 0--1 in these county-level models (not 0--100 as in the individual-level longevity models), so each coefficient is the contrast between a perfectly integrated and a perfectly segregated county; divide by ten to read it per 10-point rise in D.},
]                     %% tabularray outer close
{                     %% tabularray inner open
colspec={Q[]Q[]Q[]Q[]Q[]},
column{2,3,4,5}={}{halign=c,},
column{1}={}{halign=l,},
hline{6}={1,2,3,4,5}{solid, black, 0.05em},
}                     %% tabularray inner close
\toprule
 & Taxes & Medicaid & Health & Welfare \\ \midrule %% TinyTableHeader
D (OLS) & -193.736*** & 1.960* & -38.418 & 16.747 \\
 & (33.672) & (0.956) & (25.196) & (10.339) \\
D (IV) & -812.221*** & 0.796 & -22.033 & -10.725 \\
 & (148.076) & (5.562) & (115.648) & (60.485) \\
Num.Obs. & 7141 & 7141 & 7141 & 7141 \\
\bottomrule
\end{talltblr}
\end{table}
